\documentclass[aspectratio=169, 11pt]{beamer}

%% Paquetes esenciales
\usepackage[utf8]{inputenc}
\usepackage[spanish]{babel}
\usepackage{tikz}
\usetikzlibrary{calc, positioning, arrows.meta}
\usepackage{booktabs}
\usepackage{amsmath}
\usepackage{array}
\usepackage{hyperref}
\usepackage{graphicx}
\usepackage{colortbl} % <-- AÑADIDO PARA SOLUCIONAR ERROR DE \cellcolor

%% Tema y colores azules profesionales
\usetheme{Madrid}
\usecolortheme{dolphin}
\setbeamercolor{structure}{fg=blue!80!black}
\setbeamercolor{frametitle}{fg=white, bg=blue!80!black}
\setbeamercolor{block title}{fg=white, bg=blue!70!black}
\setbeamercolor{block body}{fg=black, bg=blue!5!white}
\setbeamercolor{normal text}{fg=black, bg=white}
\setbeamercolor{itemize item}{fg=blue!80!black}
\setbeamercolor{enumerate item}{fg=blue!80!black}
\setbeamertemplate{navigation symbols}{}
\setbeamertemplate{footline}[frame number]

%% Información del documento
\title[Modelado Multiescala: Cáncer de Mama]{Modelado multiescala de la red metabólica del cáncer de mama:\\Redes complejas y dinámica no lineal}
\author[Alejandro Guerra González]{Alejandro Guerra González}
\institute[UASLP]{
	Posgrado de Ciencias en Ingeniería Química\\
	Facultad de Ciencias Químicas\\
	Universidad Autónoma de San Luis Potosí
}
\date{\today}

\begin{document}
	
	%% Diapositiva de Título
	\begin{frame}
		\begin{figure}[htbp]
			\centering
			\includegraphics[width=0.8\textwidth]{mi_imagen.png}
			\caption{Descripción breve de la imagen}
			\label{fig:mi_imagen}
		\end{figure}
	\end{frame}
	
	%% Índice
	\begin{frame}{Índice}
		\tableofcontents
	\end{frame}
	
	\section{Introducción y Justificación}
	\begin{frame}{Introducción y Justificación}
		\begin{itemize}
			\item \textbf{Contexto global (OMS, 2020):}
			\begin{itemize}
				\item Cáncer: 1\textsuperscript{ra} causa de muerte mundial ($\sim$10 millones de defunciones)
				\item Cáncer de mama: 2.26 millones de casos, 685,000 defunciones
			\end{itemize}
			
			\vspace{0.3cm}
			\item \textbf{Naturaleza del problema:}
			\begin{itemize}
				\item Enfermedad sistémica con reprogramación metabólica profunda
				\item Heterogeneidad intratumoral y alta capacidad de adaptación
				\item Resistencia a perturbaciones terapéuticas
			\end{itemize}
			
			\vspace{0.3cm}
			\item \textbf{Necesidad de enfoques integradores:}
			\begin{itemize}
				\item Superar el reduccionismo tradicional
				\item Capturar interacciones no lineales y organización modular
				\item Dinámica multiescala del sistema
			\end{itemize}
		\end{itemize}
	\end{frame}
	
	\begin{frame}{Redes Complejas en Oncología}
		\begin{columns}
			\begin{column}{0.48\textwidth}
				\textbf{Aportes principales:}
				\begin{itemize}
					\item Caracterización de subtipos moleculares (DCloc, DCglob)
					\item Identificación de hubs proteicos (FPR2, CXCR4, GPR37)
					\item Visualización de comunicación paracrina
					\item Análisis de miRNAs (miR-200, miR-199)
				\end{itemize}
			\end{column}
			\begin{column}{0.48\textwidth}
				\textbf{Fronteras tecnológicas:}
				\begin{itemize}
					\item Gemelos Digitales multicapa
					\item Redes neuronales hipercomplejas
					\item Mapas de atención para mamografías
					\item Detección temprana mejorada
				\end{itemize}
			\end{column}
		\end{columns}
		
		\vspace{0.3cm}
		\centering
		\footnotesize \textit{El cáncer como sistema adaptativo complejo con complejidad algorítmica intratable}
	\end{frame}
	
	\section{Conexión con Modelos Dinámicos}
	\begin{frame}{Conexión con Modelos Dinámicos}
		El formalismo de redes complejas permite extender el análisis estructural hacia modelos dinámicos:
		
		\begin{equation}
			\dot{x}_i = f(x_i) + \sum_{j} A_{ij} \, g(x_i, x_j)
		\end{equation}
		
		donde:
		\begin{itemize}
			\item $x_i$: Actividad del nodo $i$ (metabolito, enzima o vía)
			\item $A_{ij}$: Matriz de adyacencia que codifica la estructura de interacciones
			\item $f(x_i)$: Dinámica interna de cada comunidad metabólica
			\item $g(x_i, x_j)$: Fuerza de interacción acoplada entre módulos
		\end{itemize}
		
		\vspace{0.3cm}
		\textbf{Paradoja tumoral:}
		\begin{itemize}
			\item \textit{Robustez estructural} $\neq$ \textit{Estabilidad dinámica}
			\item Resistente a perturbaciones aleatorias pero sensible a ataques dirigidos
			\item Coexistencia de redundancia y sensibilidad paramétrica
		\end{itemize}
	\end{frame}
	
	\section{Marco Teórico}
	\begin{frame}{Marco Teórico I: Redes Complejas}
		\begin{itemize}
			\item \textbf{Definición:} Grafo con propiedades emergentes no triviales
			\begin{itemize}
				\item Distribuciones de grado no aleatorias
				\item Modularidad y presencia de hubs
				\item Topologías \textit{scale-free} (libres de escala)
			\end{itemize}
			
			\vspace{0.2cm}
			\item \textbf{Redes metabólicas:}
			\begin{itemize}
				\item Nodos: Metabolitos o reacciones
				\item Aristas: Transformaciones bioquímicas
				\item Robustez frente a fallos aleatorios
				\item Vulnerabilidad a ataques dirigidos
			\end{itemize}
		\end{itemize}
		
		\vspace{0.3cm}
		\centering
		% AQUÍ IRÍA LA IMAGEN SI LA TUVIERAS: \includegraphics[width=0.6\textwidth]{network_topology.png}
		\framebox[0.6\textwidth]{\centering [Espacio para imagen: network\_topology]}
		\footnotesize \textit{Topología de redes: Del "hairball" al orden modular}
	\end{frame}
	
	\begin{frame}{Marco Teórico II: Modularidad y Comunidades}
		\begin{itemize}
			\item \textbf{Algoritmo de Louvain:}
			\begin{itemize}
				\item Detección de comunidades metabólicas
				\item Identificación de módulos funcionales altamente interconectados
				\item Escala mesoscópica natural para reducción dinámica
			\end{itemize}
			
			\vspace{0.2cm}
			\item \textbf{Robustez:}
			\begin{itemize}
				\item \textit{Estructural}: Mantenimiento de conectividad global
				\item \textit{Dinámica}: Persistencia de estados estacionarios
				\item Asociada a capacidad adaptativa y resistencia terapéutica
			\end{itemize}
		\end{itemize}
		
		\vspace{0.3cm}
		\centering
		% AQUÍ IRÍA LA IMAGEN SI LA TUVIERAS: \includegraphics[width=0.7\textwidth]{louvain_algorithm.png}
		\framebox[0.7\textwidth]{\centering [Espacio para imagen: louvain\_algorithm]}
		\footnotesize \textit{Detección de comunidades: Reducción de complejidad (D×10)}
	\end{frame}
	
	\begin{frame}{Marco Teórico III: Metabolismo del Cáncer de Mama}
		\begin{columns}
			\begin{column}{0.48\textwidth}
				\textbf{Metabolitos centrales:}
				\begin{itemize}
					\item Glucosa
					\item Lactato
					\item Glutamina
					\item NADH/NAD\textsuperscript{+}
					\item ATP
				\end{itemize}
			\end{column}
			\begin{column}{0.48\textwidth}
				\textbf{Rutas clave:}
				\begin{itemize}
					\item Glucólisis aeróbica (Warburg)
					\item Metabolismo de glutamina
					\item Metabolismo de lípidos
					\item PI3K/Akt/mTOR
					\item HIF-1$\alpha$
				\end{itemize}
			\end{column}
		\end{columns}
		
		\vspace{0.3cm}
		\centering
		% AQUÍ IRÍA LA IMAGEN SI LA TUVIERAS: \includegraphics[width=0.8\textwidth]{metabolic_pathways.png}
		\framebox[0.8\textwidth]{\centering [Espacio para imagen: metabolic\_pathways]}
		\footnotesize \textit{Integración de metabolismo, señalización y microambiente}
	\end{frame}
	
	\section{Hipótesis y Objetivos}
	\begin{frame}{Hipótesis de Trabajo}
		\begin{quote}
			\emph{La reducción multiescala de la red metabólica del cáncer de mama a nivel de comunidades permite identificar sistemas dinámicos efectivos que exhiben robustez estructural, complejidad dinámica e inestabilidades de tipo Turing responsables de la heterogeneidad metabólica espacial.}
		\end{quote}
		
		\vspace{0.5cm}
		\centering
		\textbf{Elementos clave:}
		\begin{itemize}
			\item Reducción multiescala por comunidades
			\item Sistemas dinámicos efectivos
			\item Robustez estructural y complejidad dinámica
			\item Patrones de Turing y heterogeneidad espacial
		\end{itemize}
	\end{frame}
	
	\begin{frame}{Objetivos}
		\textbf{Objetivo General:}
		\begin{itemize}
			\item Desarrollar y analizar un modelo dinámico multiescala de la red metabólica del cáncer de mama, integrando redes complejas, reducción por comunidades y modelos de reacción–difusión para estudiar robustez y complejidad dinámica.
		\end{itemize}
		
		\vspace{0.3cm}
		\textbf{Objetivos Específicos:}
		\begin{enumerate}
			\item Reconstruir y caracterizar la red metabólica como red compleja
			\item Identificar comunidades metabólicas funcionales (Louvain)
			\item Derivar sistema de EDO reducido (separación de escalas + renormalización)
			\item Analizar estabilidad, robustez y sensibilidad paramétrica
			\item Explorar patrones espaciales de tipo Turing
		\end{enumerate}
	\end{frame}
	
	\section{Metodología}
	\begin{frame}{Metodología: Las 5 Fases del Proyecto}
		\begin{columns}
			\begin{column}{0.48\textwidth}
				\textbf{Fase 1:} Reconstrucción y análisis topológico
				\begin{itemize}
					\item KEGG, STRING
					\item Grado, centralidad, clustering
				\end{itemize}
				
				\vspace{0.2cm}
				\textbf{Fase 2:} Detección de comunidades
				\begin{itemize}
					\item Algoritmo de Louvain
					\item Módulos funcionales
				\end{itemize}
				
				\vspace{0.2cm}
				\textbf{Fase 3:} Reducción dinámica
				\begin{itemize}
					\item Separación de escalas
					\item Renormalización
				\end{itemize}
			\end{column}
			\begin{column}{0.48\textwidth}
				\textbf{Fase 4:} Análisis dinámico
				\begin{itemize}
					\item Estabilidad
					\item Índices de Sobol
					\item Robustez
				\end{itemize}
				
				\vspace{0.2cm}
				\textbf{Fase 5:} Patrones espaciales
				\begin{itemize}
					\item Reacción-difusión
					\item Inestabilidades de Turing
					\item Atlas de Turing
				\end{itemize}
			\end{column}
		\end{columns}
		
		\vspace{0.3cm}
		\centering
		% AQUÍ IRÍA LA IMAGEN SI LA TUVIERAS: \includegraphics[width=0.9\textwidth]{methodology_phases.png}
		\framebox[0.9\textwidth]{\centering [Espacio para imagen: methodology\_phases]}
		\footnotesize \textit{Marco de trabajo estandarizado en Python}
	\end{frame}
	
	\begin{frame}{Fase 1-2: De la Red Completa a las Comunidades}
		\centering
		% AQUÍ IRÍA LA IMAGEN SI LA TUVIERAS: \includegraphics[width=0.85\textwidth]{reduction_pipeline.png}
		\framebox[0.85\textwidth]{\centering [Espacio para imagen: reduction\_pipeline]}
		
		\vspace{0.3cm}
		\textbf{Proceso de reducción:}
		\begin{enumerate}
			\item \textbf{Caos:} Red metabólica global (KEGG)
			\item \textbf{Orden:} Detección de módulos funcionales (Louvain)
			\item \textbf{Síntesis:} Una ecuación representativa por comunidad
		\end{enumerate}
		
		\vspace{0.2cm}
		\footnotesize \textit{El metabolismo tumoral actúa a través de unidades coherentes, no gen por gen}
	\end{frame}
	
	\begin{frame}{Fase 3-5: Traducción Matemática y Espacial}
		\begin{itemize}
			\item \textbf{Sistema de EDOs efectivo:}
			\begin{equation}
				\dot{x}_i = f(x_i) + \sum_{j} A_{ij} \, g(x_i, x_j)
			\end{equation}
			
			\vspace{0.2cm}
			\item \textbf{Incorporación espacial (Reacción-Difusión):}
			\begin{equation}
				\frac{\partial x_i}{\partial t} = f(x_i) + D_i \nabla^2 x_i + \sum_{j} A_{ij} \, g(x_i, x_j)
			\end{equation}
			donde $D_i$ es el coeficiente de difusión del metabolito $i$
		\end{itemize}
		
		\vspace{0.3cm}
		\centering
		% AQUÍ IRÍA LA IMAGEN SI LA TUVIERAS: \includegraphics[width=0.75\textwidth]{reaction_diffusion.png}
		\framebox[0.75\textwidth]{\centering [Espacio para imagen: reaction\_diffusion]}
		\footnotesize \textit{De la dinámica temporal a los patrones espaciotemporales}
	\end{frame}
	
	\begin{frame}{El Atlas de Turing Metabólico}
		\centering
		% AQUÍ IRÍA LA IMAGEN SI LA TUVIERAS: \includegraphics[width=0.9\textwidth]{turing_atlas.png}
		\framebox[0.9\textwidth]{\centering [Espacio para imagen: turing\_atlas]}
		
		\vspace{0.3cm}
		\textbf{Concepto:}
		\begin{itemize}
			\item La heterogeneidad intratumoral no es aleatoria
			\item Es un patrón físico-matemático predecible
			\item Mapeo de parámetros metabólicos $\rightarrow$ patrones espaciales
			\item Conexión: Topología $\rightarrow$ EDOs $\rightarrow$ Difusión $\rightarrow$ Patrones
		\end{itemize}
	\end{frame}
	
	\section{Cronograma}
	\begin{frame}{Calendario de Actividades (24 meses)}
		\renewcommand{\arraystretch}{1.3}
		\begin{tabular}{lcccc}
			\toprule
			\textbf{Actividad} & \textbf{1-6} & \textbf{7-12} & \textbf{13-18} & \textbf{19-24} \\
			\midrule
			Revisión bibliográfica & \cellcolor{blue!20}X & & & \\
			Reconstrucción de la red & \cellcolor{blue!20}X & \cellcolor{blue!20}X & & \\
			Análisis topológico y comunidades & & \cellcolor{blue!20}X & & \\
			Modelo dinámico y reducción & & & \cellcolor{blue!20}X & \cellcolor{blue!20}X \\
			Análisis de robustez y sensibilidad & & & \cellcolor{blue!20}X & \\
			Modelos de reacción-difusión & & & & \cellcolor{blue!20}X \\
			Atlas de Turing & & & & \cellcolor{blue!20}X \\
			Redacción final & & & & \cellcolor{blue!20}X \\
			\bottomrule
		\end{tabular}
		
		\vspace{0.3cm}
		\centering
		\footnotesize \textit{Cronograma estructurado en 5 fases secuenciales}
	\end{frame}
	
	\section{Resultados e Impacto}
	\begin{frame}{Resultados Esperados}
		\begin{itemize}
			\item \textbf{Caracterización topológica completa:}
			\begin{itemize}
				\item Métricas de red (grado, centralidad, clustering, modularidad)
				\item Identificación de hubs metabólicos
			\end{itemize}
			
			\vspace{0.2cm}
			\item \textbf{Modelo dinámico reducido:}
			\begin{itemize}
				\item Sistema de EDOs interpretable
				\item Una ecuación por comunidad metabólica
			\end{itemize}
			
			\vspace{0.2cm}
			\item \textbf{Cuantificación de robustez:}
			\begin{itemize}
				\item Índices de sensibilidad de Sobol
				\item Análisis de estabilidad
			\end{itemize}
			
			\vspace{0.2cm}
			\item \textbf{Atlas de Turing metabólico:}
			\begin{itemize}
				\item Mapeo parámetros $\rightarrow$ patrones espaciales
				\item Predicción de heterogeneidad intratumoral
			\end{itemize}
		\end{itemize}
	\end{frame}
	
	\begin{frame}{Impacto Esperado}
		\begin{columns}
			\begin{column}{0.32\textwidth}
				\centering
				\textbf{Entendimiento Sistémico}
				\vspace{0.2cm}
				% AQUÍ IRÍA LA IMAGEN SI LA TUVIERAS: \includegraphics[width=0.8\textwidth]{impact_systemic.png}
				\framebox[0.8\textwidth]{\centering [img]}
				\footnotesize Descifrar reglas ocultas del metabolismo
			\end{column}
			\begin{column}{0.32\textwidth}
				\centering
				\textbf{Vulnerabilidades Terapéuticas}
				\vspace{0.2cm}
				% AQUÍ IRÍA LA IMAGEN SI LA TUVIERAS: \includegraphics[width=0.8\textwidth]{impact_therapeutic.png}
				\framebox[0.8\textwidth]{\centering [img]}
				\footnotesize Modelo \textit{in silico} para simulación
			\end{column}
			\begin{column}{0.32\textwidth}
				\centering
				\textbf{Gemelos Digitales}
				\vspace{0.2cm}
				% AQUÍ IRÍA LA IMAGEN SI LA TUVIERAS: \includegraphics[width=0.8\textwidth]{impact_digital.png}
				\framebox[0.8\textwidth]{\centering [img]}
				\footnotesize Hipótesis comprobables
			\end{column}
		\end{columns}
		
		\vspace{0.3cm}
		\textbf{Contribuciones:}
		\begin{itemize}
			\item Conexión estructura-dinámica-espacio
			\item Hipótesis testables sobre heterogeneidad metabólica
			\item Identificación de vulnerabilidades terapéuticas basadas en topología
		\end{itemize}
	\end{frame}
	
	%% Diapositiva final
	\begin{frame}
		\centering
		\Huge \textbf{¡Gracias por su atención!}
		\vspace{1cm}
		
		\Large \textbf{Preguntas y comentarios}
		
		\vspace{1cm}
		\normalsize
		Alejandro Guerra González\\
		Posgrado de Ciencias en Ingeniería Química\\
		Facultad de Ciencias Químicas, UASLP\\
		\vspace{0.3cm}
		Asesor: Dr. Erik César Herrera Hernández
	\end{frame}
	
\end{document}